\chapter{Matrix Decompositions: Eigenvalues, SVD, and PCA}
\label{chap:ch04}

% Chapter goal: break a matrix into simple pieces (directions and
% stretch factors) and use them to understand what the matrix does,
% from gradient descent to PCA.

A general matrix mixes its inputs in a way that is hard to read off from
its entries. The matrix of \Cref{ex:ch03-matvec} stretched, turned, and
skewed the unit square all at once (\Cref{fig:ch03-transforms}). This
chapter finds the special directions in which a matrix acts simply, and
uses them to take the matrix apart. The pieces explain why gradient
descent slows down in narrow valleys, how far a deep network's signals
can grow, and how to compress a dataset to its most important directions.

\section{Eigenvalues and Eigenvectors}
\label{sec:ch04-eigen}

\subsection{Directions a Matrix Does Not Turn}
\label{sec:ch04-eigen-def}

Apply the symmetric matrix
\[
  A = \begin{bmatrix} 2 & 1 \\ 1 & 2 \end{bmatrix}
\]
to a few vectors. The basis vector $\vect{e}_1 = (1, 0)$ goes to
$(2, 1)$, which points in a different direction. So does
$\vect{e}_2 = (0, 1)$, which goes to $(1, 2)$. But two directions are
special:
\[
  A \begin{bmatrix} 1 \\ 1 \end{bmatrix}
  = \begin{bmatrix} 3 \\ 3 \end{bmatrix}
  = 3 \begin{bmatrix} 1 \\ 1 \end{bmatrix},
  \qquad
  A \begin{bmatrix} -1 \\ 1 \end{bmatrix}
  = \begin{bmatrix} -1 \\ 1 \end{bmatrix}
  = 1 \begin{bmatrix} -1 \\ 1 \end{bmatrix}.
\]
Along the diagonal $(1, 1)$, the matrix only stretches, by a factor of
$3$. Along the other diagonal $(-1, 1)$, it does nothing at all. In these
two directions, multiplying by a matrix reduces to multiplying by a
number.

\begin{definition}[Eigenvalue and eigenvector]
  \label{def:eigen}
  Let $A$ be an $n \times n$ matrix. A nonzero vector $\vect{v}$ is an
  \emph{eigenvector} of $A$ with \emph{eigenvalue} $\lambda$ if
  \[
    A\vect{v} = \lambda \vect{v} .
  \]
\end{definition}

The prefix \emph{eigen} is German for ``own'': an eigenvector is a
direction that belongs to the matrix. A few remarks on the definition:
\begin{itemize}
  \item Only square matrices have eigenvectors, since $A\vect{v}$ must
    live in the same space as $\vect{v}$ for the two to be compared.
  \item The vector $\vect{0}$ is excluded, because $A\vect{0} =
    \lambda\vect{0}$ holds for every $\lambda$ and says nothing. The
    eigenvalue, however, may be $0$.
  \item Eigenvectors are really directions. If $\vect{v}$ is an
    eigenvector, so is $c\vect{v}$ for any $c \ne 0$, with the same
    eigenvalue, since $A(c\vect{v}) = cA\vect{v} = \lambda (c\vect{v})$.
    We usually report the unit eigenvector, here
    $\tfrac{1}{\sqrt2}(1, 1)$ and $\tfrac{1}{\sqrt2}(-1, 1)$, and its sign
    is arbitrary.
\end{itemize}

\subsection{Geometric Interpretation}
\label{sec:ch04-eigen-geometry}

\Cref{fig:ch04-eigen} shows what $A$ does to the unit circle. Every
vector on the circle is sent somewhere on an ellipse. Most vectors are
both stretched and turned, as $\vect{e}_1$ is. The eigenvectors are the
two directions that are not turned, and they are the axes of the ellipse:
the long axis, of half-length $3$, lies along $(1, 1)$, and the short
axis, of half-length $1$, along $(-1, 1)$.

\begin{figure}[ht]
  \centering
  \begin{tikzpicture}[scale=1.3, >=stealth]
    \draw[gray!25] (-2.6,-2.6) grid (2.6,2.6);
    \draw[gray!60] (-2.6,0) -- (2.6,0);
    \draw[gray!60] (0,-2.6) -- (0,2.6);
    % image of the unit circle
    \fill[blue!10, rotate=45] (0,0) ellipse [x radius=3, y radius=1];
    \draw[blue!50!black, rotate=45] (0,0) ellipse [x radius=3, y radius=1];
    \draw[dashed, gray!70!black] (0,0) circle (1);
    % eigen directions
    \draw[green!45!black, dotted, thick] (-2.5,-2.5) -- (2.5,2.5);
    \draw[green!45!black, dotted, thick] (1.6,-1.6) -- (-1.6,1.6);
    % a vector that turns
    \draw[->, thick, gray!70!black] (0,0) -- (1,0)
      node[below right, font=\small, fill=white, inner sep=1pt]
      {$\vect{e}_1$};
    \draw[->, very thick, orange!85!black] (0,0) -- (2,1)
      node[right, font=\small] {$A\vect{e}_1 = (2, 1)$};
    % eigenvectors
    \draw[->, very thick, green!45!black] (0,0) -- (2.121,2.121)
      node[right=3pt, font=\small, fill=white, inner sep=1pt]
      {$A\vect{v}_1 = 3\vect{v}_1$};
    \draw[->, thick, gray!70!black] (0,0) -- (0.707,0.707);
    \node[gray!30!black, font=\small] at (0.3,0.78) {$\vect{v}_1$};
    \draw[->, very thick, green!45!black] (0,0) -- (-0.707,0.707)
      node[left=4pt, font=\small, fill=white, inner sep=1pt]
      {$A\vect{v}_2 = \vect{v}_2$};
  \end{tikzpicture}
  \caption{The matrix $A = \bigl[\begin{smallmatrix} 2 & 1 \\ 1 & 2
    \end{smallmatrix}\bigr]$ maps the unit circle (dashed) onto an ellipse
    (blue). A typical vector such as $\vect{e}_1$ is turned as well as
    stretched. The eigenvectors $\vect{v}_1 = \tfrac{1}{\sqrt2}(1, 1)$ and
    $\vect{v}_2 = \tfrac{1}{\sqrt2}(-1, 1)$ (green) keep their direction and
    are scaled by their eigenvalues $3$ and $1$; they point along the axes
    of the ellipse.}
  \label{fig:ch04-eigen}
\end{figure}

The eigenvalue says what happens along its eigenvector:
\begin{itemize}
  \item $\lambda > 1$: vectors in that direction are stretched;
  \item $0 < \lambda < 1$: they shrink;
  \item $\lambda < 0$: they flip to point the opposite way, and are
    stretched or shrunk by $|\lambda|$;
  \item $\lambda = 0$: they are crushed to $\vect{0}$. A matrix with a zero
    eigenvalue loses information, because $A\vect{v} = \vect{0} =
    A\vect{0}$ means two different inputs give the same output.
\end{itemize}

\paragraph{Eigenvectors of familiar transformations.}
The transformations of \Cref{sec:ch03-plane-examples} show the range of
possibilities, summarized in \Cref{tab:ch04-eigen-examples}. For the
scaling $\operatorname{diag}(2, 0.5)$, the axes themselves are the
eigenvectors, with eigenvalues $2$ and $0.5$: a diagonal matrix lists its
eigenvalues on its diagonal. The reflection $\operatorname{diag}(-1, 1)$
flips $\vect{e}_1$ ($\lambda = -1$) and leaves $\vect{e}_2$ alone
($\lambda = 1$). The projection onto the horizontal axis keeps
$\vect{e}_1$ ($\lambda = 1$) and crushes $\vect{e}_2$ ($\lambda = 0$),
which is exactly the information it loses.

Two of the examples break the pattern. The shear
$\bigl[\begin{smallmatrix} 1 & 1 \\ 0 & 1 \end{smallmatrix}\bigr]$ tips every
vector except those along $\vect{e}_1$, so it has only one eigenvector
direction, with $\lambda = 1$; there are not enough eigenvectors to
describe the whole plane. A rotation by any angle other than $0^\circ$ or
$180^\circ$ turns \emph{every} vector, so it has no eigenvectors at all.
Neither problem can occur for symmetric matrices, as
\Cref{sec:ch04-symmetric} explains.

\begin{table}[ht]
  \centering
  \caption{Eigenvalues and eigenvectors of some $2 \times 2$ matrices.}
  \label{tab:ch04-eigen-examples}
  \begin{tabular}{@{}llll@{}}
    \toprule
    Transformation & Matrix & Eigenvalues & Eigenvectors \\
    \midrule
    Scaling & $\operatorname{diag}(2, 0.5)$ & $2$, $0.5$
      & $\vect{e}_1$, $\vect{e}_2$ \\
    Reflection & $\operatorname{diag}(-1, 1)$ & $-1$, $1$
      & $\vect{e}_1$, $\vect{e}_2$ \\
    Projection & $\operatorname{diag}(1, 0)$ & $1$, $0$
      & $\vect{e}_1$, $\vect{e}_2$ \\
    Shear & $\bigl[\begin{smallmatrix} 1 & 1 \\ 0 & 1
      \end{smallmatrix}\bigr]$ & $1$ (twice) & $\vect{e}_1$ only \\
    Rotation by $90^\circ$ & $\bigl[\begin{smallmatrix} 0 & -1 \\ 1 & 0
      \end{smallmatrix}\bigr]$ & none & none \\
    Symmetric example & $\bigl[\begin{smallmatrix} 2 & 1 \\ 1 & 2
      \end{smallmatrix}\bigr]$ & $3$, $1$ & $(1, 1)$, $(-1, 1)$ \\
    \bottomrule
  \end{tabular}
\end{table}

\paragraph{Repeated multiplication.}
Eigenvectors make repeated multiplication easy. If $A\vect{v} =
\lambda\vect{v}$, then $A^2\vect{v} = A(\lambda\vect{v}) =
\lambda^2\vect{v}$, and in general
\[
  A^k \vect{v} = \lambda^k \vect{v} .
\]
For any other vector, write it as a combination of eigenvectors and apply
this to each piece. With the matrix $A$ above, $\vect{e}_1 = \tfrac12 (1,
1) + \tfrac12 (1, -1)$, so
\[
  A^k \vect{e}_1
  = \tfrac12\, 3^k \begin{bmatrix} 1 \\ 1 \end{bmatrix}
  + \tfrac12\, 1^k \begin{bmatrix} 1 \\ -1 \end{bmatrix}
  = \frac12 \begin{bmatrix} 3^k + 1 \\ 3^k - 1 \end{bmatrix}.
\]
The iterates $(1, 0), (2, 1), (5, 4), (14, 13), (41, 40), \dots$ swing
toward the direction $(1, 1)$, because the component along the
eigenvector with the largest $|\lambda|$ grows fastest and soon dominates.
This is the idea behind \emph{power iteration}: multiply any starting
vector by $A$ over and over, normalizing as you go, and it lines up with
the \emph{dominant} eigenvector. More generally, the size of $A^k\vect{x}$
over many steps is governed by the eigenvalues: components with
$|\lambda| > 1$ explode, and components with $|\lambda| < 1$ die out.

\begin{aibox}
  A recurrent neural network applies the same weight matrix $W$ at every
  time step, so signals and gradients passed back through $k$ steps are
  multiplied by something like $W^k$. If the largest $|\lambda|$ is above
  $1$, they explode; if it is below $1$, they vanish, and the network
  cannot learn long-range dependencies. This \emph{exploding and vanishing
  gradient problem} motivated LSTMs, careful initialization schemes, and
  gradient clipping. Google's original PageRank algorithm ranked web pages
  by the dominant eigenvector of the web's link matrix, computed with
  power iteration.
\end{aibox}

\paragraph{Eigenvalues and gradient descent.}
In \Cref{sec:ch02-gd-multi}, gradient descent on $x_1^2 + 3x_2^2$ needed
a step size below $1/3$, because each coordinate was multiplied by its
own factor, $1 - 2\eta$ or $1 - 6\eta$, at every step. Eigenvalues explain
what happens when the valley is not lined up with the axes. Take
$f(\vect{x}) = \vect{x}^\top A \vect{x} = 2x_1^2 + 2x_1 x_2 + 2x_2^2$, whose
gradient is $\nabla f(\vect{x}) = 2A\vect{x}$. One step of gradient
descent is
\[
  \vect{x} \leftarrow \vect{x} - 2\eta A\vect{x} = (I - 2\eta A)\,\vect{x} ,
\]
a multiplication by the matrix $I - 2\eta A$. Its eigenvectors are those
of $A$, with eigenvalues $1 - 2\eta\lambda$, that is, $1 - 6\eta$ along
$(1, 1)$ and $1 - 2\eta$ along $(-1, 1)$. This is the same pair of factors
as before, now acting along the eigenvectors instead of the coordinate
axes. The level curves of $f$ are ellipses whose axes are the
eigenvectors, the steep direction has the large eigenvalue, and the step
size must again satisfy $\eta < 1/3$. In general, the largest eigenvalue
of the matrix of second derivatives sets the largest stable learning
rate, and the ratio of the largest to the smallest eigenvalue measures
how slowly gradient descent crawls along the flat directions.

\subsection{Finding Eigenvalues}
\label{sec:ch04-eigen-compute}

Rewrite $A\vect{v} = \lambda\vect{v}$ as
\[
  (A - \lambda I)\,\vect{v} = \vect{0} .
\]
We are looking for a nonzero $\vect{v}$ that the matrix $A - \lambda I$
sends to $\vect{0}$. As we saw with the projection, this happens exactly
when $A - \lambda I$ flattens the space onto something lower-dimensional
and loses information, which by \Cref{prop:ch03-invertible} is exactly
when its determinant is zero. Recall from \Cref{sec:ch03-det} that the
determinant is the factor by which a matrix scales areas, and for a
$2 \times 2$ matrix it is $ad - bc$. For example, $\det A = 4 - 1 = 3$,
and the ellipse in \Cref{fig:ch04-eigen} has area $3\pi$, three times
that of the unit circle. The eigenvalues are therefore the solutions of
the \emph{characteristic equation}
\[
  \det(A - \lambda I) = 0 .
\]

\begin{example}
  \label{ex:ch04-eigen-2x2}
  For $A = \bigl[\begin{smallmatrix} 2 & 1 \\ 1 & 2 \end{smallmatrix}\bigr]$,
  \[
    \det(A - \lambda I)
    = \det \begin{bmatrix} 2 - \lambda & 1 \\ 1 & 2 - \lambda \end{bmatrix}
    = (2 - \lambda)^2 - 1
    = \lambda^2 - 4\lambda + 3
    = (\lambda - 3)(\lambda - 1),
  \]
  so the eigenvalues are $\lambda = 3$ and $\lambda = 1$. To find the
  eigenvector for $\lambda = 3$, solve $(A - 3I)\vect{v} = \vect{0}$:
  \[
    \begin{bmatrix} -1 & 1 \\ 1 & -1 \end{bmatrix}
    \begin{bmatrix} v_1 \\ v_2 \end{bmatrix} = \vect{0}
    \quad\Longrightarrow\quad v_1 = v_2 ,
  \]
  which gives the direction $(1, 1)$. The same steps for $\lambda = 1$
  give $v_1 = -v_2$, the direction $(-1, 1)$.
\end{example}

For any $2 \times 2$ matrix, the characteristic equation is
$\lambda^2 - (a + d)\lambda + (ad - bc) = 0$. The coefficient $a + d$ is
the \emph{trace}, the sum of the diagonal entries, so
\[
  \lambda_1 + \lambda_2 = \operatorname{tr} A, \qquad
  \lambda_1 \lambda_2 = \det A .
\]
Both facts hold for $n \times n$ matrices too: the eigenvalues add up to
the trace and multiply to the determinant. They make quick checks: for
our $A$, $3 + 1 = 4 = 2 + 2$ and $3 \cdot 1 = 3 = \det A$.

For an $n \times n$ matrix, $\det(A - \lambda I)$ is a polynomial of
degree $n$, so there are at most $n$ eigenvalues. Solving that
polynomial is fine by hand for $n = 2$ and
hopeless for the matrices of machine learning. Numerical libraries use
iterative algorithms instead, such as the QR algorithm, which never form
the polynomial at all. In practice, you compute eigenvalues with a library
call, as in \Cref{sec:ch04-eigen-code}.

\subsection{Symmetric Matrices}
\label{sec:ch04-symmetric}

The shear and the rotation show that a general matrix may have too few
eigenvectors, or none at all. The matrices we meet most often in
machine learning avoid both problems because they are \emph{symmetric}:
$A^\top = A$. Covariance matrices of data, Gram matrices $EE^\top$ from
\Cref{sec:ch03-matmul-ai}, and matrices of second derivatives are all
symmetric.

\begin{theorem}[Spectral theorem]
  \label{thm:spectral}
  A symmetric $n \times n$ matrix has $n$ eigenvalues, counted with
  repetition, and a set of $n$ unit eigenvectors that are mutually
  orthogonal.
\end{theorem}

We will not prove the whole theorem, but the orthogonality is easy to see.
Suppose $A\vect{v} = \lambda\vect{v}$ and $A\vect{w} = \mu\vect{w}$ with
$\lambda \ne \mu$. Using $A^\top = A$,
\[
  \lambda\, \vect{v}^\top \vect{w}
  = (A\vect{v})^\top \vect{w}
  = \vect{v}^\top A^\top \vect{w}
  = \vect{v}^\top A \vect{w}
  = \mu\, \vect{v}^\top \vect{w} ,
\]
so $(\lambda - \mu)\, \vect{v}^\top \vect{w} = 0$, which forces
$\vect{v}^\top \vect{w} = 0$. Eigenvectors of a symmetric matrix with
different eigenvalues are perpendicular, as $(1, 1)$ and $(-1, 1)$ are in
\Cref{fig:ch04-eigen}.

Geometrically, the theorem says that every symmetric matrix acts like the
example in this section: there are $n$ perpendicular directions, and the
matrix stretches or flips the space along each of them by its eigenvalue.
The unit sphere becomes an ellipsoid whose axes are the eigenvectors. The
next section turns this picture into a factorization of the matrix
itself.

\subsection{In Code: Computing Eigenvalues}
\label{sec:ch04-eigen-code}

NumPy computes eigenvalues and eigenvectors with \texttt{np.linalg.eig}.
It returns the eigenvalues as a vector and the unit eigenvectors as the
\emph{columns} of a matrix:

\begin{lstlisting}
import numpy as np

A = np.array([[2.0, 1.0],
              [1.0, 2.0]])

eigvals, eigvecs = np.linalg.eig(A)
eigvals                    # [3., 1.]
eigvecs                    # [[0.7071, -0.7071],
                           #  [0.7071,  0.7071]]

v = eigvecs[:, 0]          # column 0, not row 0
A @ v                      # [2.1213, 2.1213]
eigvals[0] * v             # [2.1213, 2.1213]: A v = 3 v

np.allclose(A @ eigvecs, eigvecs * eigvals)   # True: all pairs at once
eigvals.sum(), eigvals.prod()   # 4.0, 3.0: trace and determinant
\end{lstlisting}

The last check uses broadcasting: \texttt{eigvecs * eigvals} multiplies
column $j$ by \texttt{eigvals[j]}, so it compares $A\vect{v}_j$ with
$\lambda_j\vect{v}_j$ for every $j$ in one line.

\begin{note}
  In NumPy 2.5 and later, \texttt{eig} prints its results with a trailing
  \texttt{+0.j}, as in \texttt{3.+0.j} instead of \texttt{3.}. The extra
  part is always zero for the matrices in this book and can be dropped
  with \texttt{.real}. For a symmetric matrix, use
  \texttt{np.linalg.eigh} instead, described next.
\end{note}

For symmetric matrices, \texttt{np.linalg.eigh} is faster and more
accurate, returns plain floating-point numbers, and sorts the eigenvalues in
increasing order. Its eigenvectors are orthonormal, so
$V^\top V = I$:

\begin{lstlisting}
eigvals, V = np.linalg.eigh(A)
eigvals                    # [1., 3.]: ascending
V                          # [[-0.7071, 0.7071],
                           #  [ 0.7071, 0.7071]]
np.allclose(V.T @ V, np.eye(2))   # True: orthonormal columns
\end{lstlisting}

The signs differ from those returned by \texttt{eig}: $-\vect{v}$ is as
good an eigenvector as $\vect{v}$, and libraries make no promise about
which one you get. Code that compares eigenvectors across runs or
libraries should allow for a sign flip.

The shear of \Cref{tab:ch04-eigen-examples}, which has too few
eigenvectors, shows up as an eigenvector that repeats:

\begin{lstlisting}
shear = np.array([[1.0, 1.0],
                  [0.0, 1.0]])
eigvals, eigvecs = np.linalg.eig(shear)
eigvals                    # [1., 1.]
eigvecs                    # [[1., -1.],
                           #  [0.,  0.]]: both columns along e1
\end{lstlisting}

Power iteration takes only a few lines. Each step multiplies by $A$ and
renormalizes; at the end, $\vect{x}^\top A \vect{x}$ estimates the
eigenvalue, since $\vect{v}^\top A \vect{v} = \lambda\, \vect{v}^\top
\vect{v} = \lambda$ for a unit eigenvector:

\begin{lstlisting}
def power_iteration(A, x0, num_iters=50):
    x = x0 / np.linalg.norm(x0)
    for _ in range(num_iters):
        x = A @ x
        x = x / np.linalg.norm(x)        # keep x at unit length
    return x @ A @ x, x

power_iteration(A, np.array([1.0, 0.0]))
# (3.0, [0.7071, 0.7071])
\end{lstlisting}

Finally, the eigenvalues of $I - 2\eta A$ predict whether gradient
descent on $f(\vect{x}) = \vect{x}^\top A \vect{x}$ converges, before we
run it:

\begin{lstlisting}
for lr in [0.1, 0.25, 0.35]:
    factors = np.linalg.eigvalsh(np.eye(2) - 2 * lr * A)
    x = np.array([-2.5, 1.2])
    for _ in range(50):
        x = x - lr * 2 * A @ x           # gradient of x^T A x is 2 A x
    print(lr, factors, x)
# 0.1   [0.4, 0.8]    [0, 0]            converges
# 0.25  [-0.5, 0.5]   [0, 0]            converges, zig-zagging
# 0.35  [-1.1, 0.3]   [-76.3, -76.3]    diverges along (1, 1)
\end{lstlisting}

The run with $\eta = 0.35$ blows up along exactly the eigenvector whose
factor $1 - 6\eta = -1.1$ has absolute value above $1$.

\section{The Singular Value Decomposition}
\label{sec:ch04-svd}

Eigenvectors describe a matrix completely only when there are enough of
them, which the spectral theorem guarantees for symmetric matrices and
nothing guarantees in general. Most matrices in machine learning are not
symmetric, and many are not even square: a data matrix has one row per
example and one column per feature. The \emph{singular value
decomposition} (SVD) extends the eigenvector picture to every matrix. We
first turn \Cref{thm:spectral} into a factorization, and then see what
changes when the matrix is not symmetric.

\subsection{From Eigenvectors to a Factorization}
\label{sec:ch04-eigendecomposition}

Let $A$ be a symmetric $n \times n$ matrix with unit eigenvectors
$\vect{q}_1, \dots, \vect{q}_n$, which by \Cref{thm:spectral} we can take
to be mutually orthogonal, and eigenvalues $\lambda_1, \dots, \lambda_n$.
Place the eigenvectors side by side as the columns of a matrix $Q$, and
the eigenvalues on the diagonal of $\Lambda =
\operatorname{diag}(\lambda_1, \dots, \lambda_n)$. By the column view of
the matrix product, the $n$ equations $A\vect{q}_i = \lambda_i \vect{q}_i$
say together that
\[
  AQ = Q\Lambda .
\]
The columns of $Q$ are orthonormal, so $Q^\top Q = I$: entry $(i, j)$ of
$Q^\top Q$ is $\vect{q}_i^\top \vect{q}_j$, which is $1$ when $i = j$ and
$0$ otherwise. A square matrix with this property is called
\emph{orthogonal}. It also satisfies $QQ^\top = I$, so $Q^\top$ undoes
$Q$ and $Q$ undoes $Q^\top$. Multiplying $AQ = Q\Lambda$ on the right by
$Q^\top$ gives the \emph{eigendecomposition}
\begin{equation}
  A = Q \Lambda Q^\top .
  \label{eq:ch04-eigendecomposition}
\end{equation}

An orthogonal matrix preserves inner products, $(Q\vect{x})^\top
(Q\vect{y}) = \vect{x}^\top Q^\top Q \vect{y} = \vect{x}^\top \vect{y}$, and
therefore lengths and angles. In the plane, that leaves only rotations
and reflections. Reading \cref{eq:ch04-eigendecomposition} from right to
left, as we read every product, multiplying by a symmetric matrix takes
three steps:
\begin{enumerate}
  \item $Q^\top$ rotates the space so that the eigenvectors line up with
    the coordinate axes ($Q^\top \vect{q}_i = \vect{e}_i$);
  \item $\Lambda$ stretches each axis by its eigenvalue;
  \item $Q$ rotates back, carrying each axis to its eigenvector.
\end{enumerate}
For the matrix of \Cref{sec:ch04-eigen},
\[
  \begin{bmatrix} 2 & 1 \\ 1 & 2 \end{bmatrix}
  =
  \underbrace{\frac{1}{\sqrt2}
    \begin{bmatrix} 1 & -1 \\ 1 & 1 \end{bmatrix}}_{Q}
  \underbrace{\begin{bmatrix} 3 & 0 \\ 0 & 1 \end{bmatrix}}_{\Lambda}
  \underbrace{\frac{1}{\sqrt2}
    \begin{bmatrix} 1 & 1 \\ -1 & 1 \end{bmatrix}}_{Q^\top} :
\]
rotate by $-45^\circ$, stretch the horizontal axis by $3$, rotate back by
$45^\circ$. Every symmetric matrix is a stretch along perpendicular axes,
seen from a rotated point of view.

The same factorization can be written as a sum of outer products, by the
outer-product view~\eqref{eq:ch03-outer} of the matrix product:
\[
  A = \lambda_1 \vect{q}_1 \vect{q}_1^\top + \dots
    + \lambda_n \vect{q}_n \vect{q}_n^\top .
\]
Each term $\vect{q}_i \vect{q}_i^\top$ is the projection onto the direction
$\vect{q}_i$, since $\vect{q}_i \vect{q}_i^\top \vect{x} = (\vect{q}_i^\top
\vect{x})\, \vect{q}_i$. A symmetric matrix splits any input into its
components along the eigenvectors and scales each by its eigenvalue.

\subsection{Definition}
\label{sec:ch04-svd-def}

For a matrix that is not symmetric, or not square, we cannot hope to find
directions that keep their direction: for an $m \times n$ matrix with
$m \ne n$, $A\vect{v}$ is not even in the same space as $\vect{v}$. The SVD
asks for less. Instead of one set of perpendicular directions, it uses
two: perpendicular \emph{input} directions $\vect{v}_i$ that $A$ sends to
perpendicular \emph{output} directions $\vect{u}_i$, scaled by a factor
$\sigma_i \ge 0$:
\begin{equation}
  A\vect{v}_i = \sigma_i \vect{u}_i .
  \label{eq:ch04-svd-pairs}
\end{equation}
Collecting these equations into matrices, exactly as we did for
eigenvectors, gives the decomposition.

\begin{theorem}[Singular value decomposition]
  \label{thm:svd}
  Every $m \times n$ matrix $A$ can be factored as
  \[
    A = U \Sigma V^\top ,
  \]
  where $U$ is an $m \times m$ orthogonal matrix, $V$ is an $n \times n$
  orthogonal matrix, and $\Sigma$ is an $m \times n$ matrix that is zero
  except for the entries $\sigma_1 \ge \sigma_2 \ge \dots \ge 0$ on its
  diagonal.
\end{theorem}

The numbers $\sigma_i$ are the \emph{singular values} of $A$, the columns
$\vect{u}_i$ of $U$ are its \emph{left singular vectors}, and the columns
$\vect{v}_i$ of $V$ are its \emph{right singular vectors}. Unlike
eigenvalues, singular values are never negative, and by convention they
are sorted from largest to smallest. Since $\Sigma$ has the same shape as
$A$, the shapes work out:
\[
  \underbrace{A}_{m \times n}
  = \underbrace{U}_{m \times m}
    \underbrace{\Sigma}_{m \times n}
    \underbrace{V^\top}_{n \times n} .
\]

\paragraph{Where the SVD comes from.}
The matrix $A^\top A$ is $n \times n$ and symmetric, since $(A^\top A)^\top
= A^\top A$. If $A = U\Sigma V^\top$, then
\[
  A^\top A = V \Sigma^\top U^\top U \Sigma V^\top = V (\Sigma^\top \Sigma)
  V^\top ,
\]
which is an eigendecomposition: the right singular vectors $\vect{v}_i$ are
the eigenvectors of $A^\top A$, and the squared singular values
$\sigma_i^2$ are its eigenvalues. This also shows how to build the SVD.
Take the orthonormal eigenvectors $\vect{v}_i$ of $A^\top A$ given by the
spectral theorem. Their eigenvalues are $\|A\vect{v}_i\|^2 =
\vect{v}_i^\top A^\top A \vect{v}_i \ge 0$, so they have square roots
$\sigma_i$. For each $\sigma_i > 0$, set $\vect{u}_i = A\vect{v}_i /
\sigma_i$. These output directions are orthonormal too:
\[
  \vect{u}_i^\top \vect{u}_j
  = \frac{\vect{v}_i^\top A^\top A \vect{v}_j}{\sigma_i \sigma_j}
  = \frac{\sigma_j^2\, \vect{v}_i^\top \vect{v}_j}{\sigma_i \sigma_j} ,
\]
which is $1$ for $i = j$ and $0$ otherwise. If there are fewer than $m$ of
them, fill out $U$ with any further orthonormal vectors; they are
multiplied by zeros in $\Sigma$ and do not affect the product.

\begin{example}
  \label{ex:ch04-svd}
  Take the matrix of \Cref{ex:ch03-matvec}, the last panel of
  \Cref{fig:ch03-transforms}:
  \[
    A = \begin{bmatrix} 2 & -1 \\ 1 & 1 \end{bmatrix}, \qquad
    A^\top A = \begin{bmatrix} 5 & -1 \\ -1 & 2 \end{bmatrix}.
  \]
  The characteristic equation of $A^\top A$ is $\lambda^2 - 7\lambda + 9 =
  0$, with roots $\lambda = (7 \pm \sqrt{13})/2$, so the singular values
  are their square roots,
  \[
    \sigma_1 = \frac{\sqrt{13} + 1}{2} \approx 2.303, \qquad
    \sigma_2 = \frac{\sqrt{13} - 1}{2} \approx 1.303 .
  \]
  The corresponding unit eigenvectors of $A^\top A$ are
  $\vect{v}_1 \approx (0.957, -0.290)$ and $\vect{v}_2 \approx (0.290,
  0.957)$, and dividing $A\vect{v}_i$ by $\sigma_i$ gives
  $\vect{u}_1 \approx (0.957, 0.290)$ and $\vect{u}_2 \approx (-0.290,
  0.957)$. Both $V$ and $U$ turn out to be rotations by
  $\theta \approx 16.8^\circ$, in opposite directions:
  \[
    A \approx
    \underbrace{\begin{bmatrix} 0.957 & -0.290 \\ 0.290 & 0.957
      \end{bmatrix}}_{U}
    \underbrace{\begin{bmatrix} 2.303 & 0 \\ 0 & 1.303
      \end{bmatrix}}_{\Sigma}
    \underbrace{\begin{bmatrix} 0.957 & -0.290 \\ 0.290 & 0.957
      \end{bmatrix}}_{V^\top} .
  \]
  Here $V^\top$ and $U$ happen to be the same rotation by $+16.8^\circ$.
  As a check, $\sigma_1 \sigma_2 = (13 - 1)/4 = 3 = \det A$.
\end{example}

If $A$ is symmetric with nonnegative eigenvalues, as in
\Cref{sec:ch04-eigendecomposition}, its SVD and its eigendecomposition
coincide, with $U = V = Q$ and $\Sigma = \Lambda$. In general they differ,
and the singular values are not the eigenvalues. The shear
$\bigl[\begin{smallmatrix} 1 & 1 \\ 0 & 1 \end{smallmatrix}\bigr]$, for
example, has the single eigenvalue $1$ but singular values $1.618$ and
$0.618$: although it keeps one direction fixed, it stretches some vectors
by more than half as much again and squeezes others.

\subsection{Geometric Interpretation}
\label{sec:ch04-svd-geometry}

Reading $A = U\Sigma V^\top$ from right to left, every matrix acts in three
steps (\Cref{fig:ch04-svd}):
\begin{enumerate}
  \item $V^\top$ rotates (or reflects) the input space so that the right
    singular vectors $\vect{v}_i$ line up with the coordinate axes;
  \item $\Sigma$ stretches axis $i$ by $\sigma_i$, and, if $m \ne n$, adds
    or drops coordinates to change the dimension from $n$ to $m$;
  \item $U$ rotates (or reflects) the output space, carrying axis $i$ to
    the left singular vector $\vect{u}_i$.
\end{enumerate}
This is the claim made at the end of \Cref{sec:ch03-plane-examples}:
however complicated a matrix looks, it is a rotation, a stretch along the
axes, and another rotation. Only the middle step changes lengths, and it
does so in the simplest possible way.

\begin{figure}[ht]
  \centering
  \begin{tikzpicture}[scale=0.62, >=stealth]
    \tikzset{
      panelaxes/.style={gray!60},
      stage/.style={->, thick, gray!70!black},
      v1/.style={->, very thick, blue!60!black},
      v2/.style={->, very thick, orange!85!black},
    }
    % (a) input: unit circle with right singular vectors
    \begin{scope}
      \draw[panelaxes] (-2.4,0) -- (2.4,0);
      \draw[panelaxes] (0,-1.8) -- (0,1.8);
      \fill[blue!10] (0,0) circle (1);
      \draw[blue!50!black] (0,0) circle (1);
      \draw[v1] (0,0) -- (0.957,-0.290) node[below right] {$\vect{v}_1$};
      \draw[v2] (0,0) -- (0.290,0.957) node[above right] {$\vect{v}_2$};
      \node[font=\small] at (0,-2.4) {(a) input};
    \end{scope}
    \draw[stage] (2.3,0.9) -- (3.7,0.9) node[midway, above] {$V^\top$};
    % (b) after V^T: singular vectors on the axes
    \begin{scope}[xshift=6cm]
      \draw[panelaxes] (-2.4,0) -- (2.4,0);
      \draw[panelaxes] (0,-1.8) -- (0,1.8);
      \fill[blue!10] (0,0) circle (1);
      \draw[blue!50!black] (0,0) circle (1);
      \draw[v1] (0,0) -- (1,0) node[below right] {$\vect{e}_1$};
      \draw[v2] (0,0) -- (0,1) node[above right] {$\vect{e}_2$};
      \node[font=\small] at (0,-2.4) {(b) rotate};
    \end{scope}
    \draw[stage] (8.3,0.9) -- (9.7,0.9) node[midway, above] {$\Sigma$};
    % (c) after Sigma: axis-aligned ellipse
    \begin{scope}[xshift=12.2cm]
      \draw[panelaxes] (-2.6,0) -- (2.6,0);
      \draw[panelaxes] (0,-1.8) -- (0,1.8);
      \fill[blue!10] (0,0) ellipse [x radius=2.303, y radius=1.303];
      \draw[blue!50!black] (0,0) ellipse [x radius=2.303, y radius=1.303];
      \draw[v1] (0,0) -- (2.303,0) node[below=2pt] {$\sigma_1 \vect{e}_1$};
      \draw[v2] (0,0) -- (0,1.303) node[above right] {$\sigma_2 \vect{e}_2$};
      \node[font=\small] at (0,-2.4) {(c) stretch};
    \end{scope}
    \draw[stage] (14.9,0.9) -- (16.3,0.9) node[midway, above] {$U$};
    % (d) after U: rotated ellipse, the image of the circle under A
    \begin{scope}[xshift=19cm]
      \draw[panelaxes] (-2.6,0) -- (2.6,0);
      \draw[panelaxes] (0,-1.8) -- (0,1.8);
      \fill[blue!10, rotate=16.845] (0,0)
        ellipse [x radius=2.303, y radius=1.303];
      \draw[blue!50!black, rotate=16.845] (0,0)
        ellipse [x radius=2.303, y radius=1.303];
      \draw[v1] (0,0) -- (2.204,0.667) node[right] {$\sigma_1 \vect{u}_1$};
      \draw[v2] (0,0) -- (-0.378,1.247) node[above] {$\sigma_2 \vect{u}_2$};
      \node[font=\small] at (0,-2.4) {(d) rotate};
    \end{scope}
  \end{tikzpicture}
  \caption{The SVD of $A = \bigl[\begin{smallmatrix} 2 & -1 \\ 1 & 1
    \end{smallmatrix}\bigr]$ as three steps. (a) The unit circle and the
    right singular vectors. (b) $V^\top$ rotates them onto the axes. (c)
    $\Sigma$ stretches the axes by $\sigma_1 \approx 2.30$ and $\sigma_2
    \approx 1.30$, turning the circle into an ellipse. (d) $U$ rotates the
    ellipse into place. The result is the image of the unit circle under
    $A$, with $A\vect{v}_i = \sigma_i \vect{u}_i$.}
  \label{fig:ch04-svd}
\end{figure}

The picture gives each part of the SVD a meaning:
\begin{itemize}
  \item Any matrix maps the unit circle (or sphere) to an ellipse (or
    ellipsoid). The singular values are the half-lengths of its axes, and
    the left singular vectors $\vect{u}_i$ point along those axes.
  \item The right singular vectors $\vect{v}_i$ are the input directions
    that land on the axes. In \Cref{fig:ch04-svd}, $\vect{v}_1$ is the
    direction that $A$ stretches the most, and $\vect{v}_2$ the one it
    stretches the least.
  \item So $\sigma_1$ is the largest factor by which $A$ can stretch any
    vector: $\|A\vect{x}\| \le \sigma_1 \|\vect{x}\|$ for every $\vect{x}$,
    with equality along $\vect{v}_1$. For a square matrix, the smallest
    singular value is the smallest stretch, and the ratio
    $\sigma_1 / \sigma_n$, the \emph{condition number}, measures how
    unevenly the matrix treats different directions.
  \item Areas scale by the product of the stretches: $|\det A| =
    \sigma_1 \sigma_2$ for a $2 \times 2$ matrix, and $|\det A| =
    \sigma_1 \cdots \sigma_n$ in general, since the rotations do not change
    volume.
\end{itemize}

\paragraph{Rank.}
Some singular values may be zero. If exactly $r$ of them are positive,
then $\Sigma$ keeps $r$ coordinates and sets the rest to zero, so every
output $A\vect{x}$ is a combination of $\vect{u}_1, \dots, \vect{u}_r$. The
number $r$ is the \emph{rank} of $A$: the number of independent
directions its outputs can take. The projection of
\Cref{sec:ch03-plane-examples} has singular values $1$ and $0$, so its
rank is $1$, and its outputs lie on a line. This is the precise version of
the information loss we noticed there: the $n - r$ input directions
$\vect{v}_{r+1}, \dots, \vect{v}_n$ are sent to $\vect{0}$, and nothing about
them can be recovered from the output.

\paragraph{A sum of rank-one pieces.}
Dropping the terms with $\sigma_i = 0$ and using the outer-product view
once more,
\begin{equation}
  A = \sigma_1 \vect{u}_1 \vect{v}_1^\top + \sigma_2 \vect{u}_2
  \vect{v}_2^\top + \dots + \sigma_r \vect{u}_r \vect{v}_r^\top .
  \label{eq:ch04-svd-sum}
\end{equation}
Each term $\vect{u}_i \vect{v}_i^\top$ is a matrix of rank one: it measures
how far the input reaches along $\vect{v}_i$ and sends the answer out along
$\vect{u}_i$. The SVD writes any matrix as a sum of these simplest
possible pieces, sorted by importance, with $\sigma_i$ as the weight of
piece $i$.

\subsection{Keeping the Largest Singular Values}
\label{sec:ch04-low-rank}

The sum~\eqref{eq:ch04-svd-sum} suggests a way to compress a matrix: keep
only the first $k$ terms,
\[
  A_k = \sigma_1 \vect{u}_1 \vect{v}_1^\top + \dots
    + \sigma_k \vect{u}_k \vect{v}_k^\top .
\]
This is the \emph{truncated SVD}, a matrix of rank $k$. Storing it takes
$k$ singular values and $k$ pairs of vectors, $k(m + n + 1)$ numbers
instead of the $mn$ entries of $A$. To measure how much is lost, we use
the \emph{Frobenius norm} $\|A\|_F = \sqrt{\sum_{i,j} a_{ij}^2}$, the norm
of the matrix viewed as one long vector. Because $U$ and $V$ preserve
lengths, it equals $\sqrt{\sigma_1^2 + \dots + \sigma_r^2}$.

\begin{theorem}[Eckart--Young]
  \label{thm:eckart-young}
  Among all $m \times n$ matrices of rank at most $k$, the truncated SVD
  $A_k$ is the closest to $A$ in the Frobenius norm, and
  \[
    \|A - A_k\|_F = \sqrt{\sigma_{k+1}^2 + \dots + \sigma_r^2} .
  \]
\end{theorem}

The error is made of exactly the singular values we threw away. When the
singular values fall off quickly, a small $k$ captures almost all of the
matrix. Real data often looks like this: a few strong patterns plus a
little noise. \Cref{fig:ch04-singular-values} shows the singular values of
a $100 \times 50$ matrix built from three random patterns plus small
random noise. Three singular values stand out, and the other $47$ are
small and nearly equal, the signature of noise. Keeping $k = 3$ terms
stores $453$ numbers instead of $5000$ and still reproduces the matrix to
within about $5\%$. It even lands closer to the noise-free matrix than the
noisy original does, because the discarded terms were mostly noise. The
next section applies the same idea to a dataset, where the top singular
vectors become its principal components.

\begin{figure}[ht]
  \centering
  \begin{tikzpicture}
    \begin{axis}[
      width=0.62\textwidth, height=0.36\textwidth,
      xmin=0, xmax=16, ymin=0, ymax=88,
      xlabel={index $i$}, ylabel={$\sigma_i$},
      xtick={1,3,5,10,15}, ytick={0,20,40,60,80},
      tick label style={font=\footnotesize},
      label style={font=\small},
      axis lines=left,
    ]
      \addplot[ycomb, very thick, blue!60!black, mark=*,
        mark size=1.8pt] coordinates {
        (1,79.17) (2,66.29) (3,63.85)};
      \addplot[ycomb, thick, gray!70, mark=*, mark size=1.5pt]
        coordinates {
        (4,1.61) (5,1.58) (6,1.50) (7,1.47) (8,1.39) (9,1.38) (10,1.33)
        (11,1.31) (12,1.28) (13,1.26) (14,1.23) (15,1.20)};
    \end{axis}
  \end{tikzpicture}
  \caption{The first $15$ singular values of a $100 \times 50$ matrix made
    of three random patterns (a rank-$3$ matrix) plus small noise. The
    three large singular values (blue) carry the patterns; the rest (gray)
    are noise.}
  \label{fig:ch04-singular-values}
\end{figure}

\begin{aibox}
  Fine-tuning a large language model normally updates every entry of
  every weight matrix. \emph{LoRA} (low-rank adaptation) instead freezes
  each weight matrix $W$ and learns a low-rank update, $W + BA$, where $B$
  is $d \times r$ and $A$ is $r \times d$ with a small rank $r$. For
  $d = 4096$ and $r = 8$, the update has $2 \cdot 4096 \cdot 8 = 65{,}536$
  parameters instead of $4096^2 \approx 16.8$ million, about $0.4\%$. The
  method works because the changes needed to adapt a pretrained model to
  a new task turn out to have most of their weight in a few singular
  directions. The same reasoning lets truncated SVDs compress embedding
  tables and the weight matrices of trained networks.
\end{aibox}

\subsection{In Code: Computing the SVD}
\label{sec:ch04-svd-code}

\texttt{np.linalg.svd} returns three arrays. Two details trip up almost
everyone: the singular values come back as a one-dimensional array, not
as the matrix $\Sigma$, and the third array is $V^\top$, not $V$.

\begin{lstlisting}
A = np.array([[2.0, -1.0],
              [1.0,  1.0]])

U, S, Vt = np.linalg.svd(A)
S                          # [2.3028, 1.3028]: sorted, largest first
U                          # [[-0.9571, -0.2898],
                           #  [-0.2898,  0.9571]]
Vt                         # [[-0.9571,  0.2898],   rows are v1, v2
                           #  [ 0.2898,  0.9571]]

np.allclose(U @ np.diag(S) @ Vt, A)        # True: A = U Sigma V^T
np.allclose(A @ Vt[0], S[0] * U[:, 0])     # True: A v1 = sigma1 u1
\end{lstlisting}

Compared with \Cref{ex:ch04-svd}, NumPy has flipped the signs of both
$\vect{u}_1$ and $\vect{v}_1$. Flipping a matching pair leaves
$\sigma_1 \vect{u}_1 \vect{v}_1^\top$ unchanged, so both answers are
correct. Because \texttt{Vt} holds $V^\top$, the right singular vectors
are its \emph{rows}, while the left singular vectors are the
\emph{columns} of \texttt{U}.

\begin{warning}
  Writing \texttt{U @ S @ Vt} fails or, worse, silently computes the wrong
  thing, because \texttt{S} is a vector. Use \texttt{U @ np.diag(S) @ Vt},
  or the broadcasting form \texttt{U * S @ Vt}, which scales column $i$ of
  \texttt{U} by \texttt{S[i]}. And remember that the right singular
  vectors are \texttt{Vt[i]}, not \texttt{Vt[:, i]}.
\end{warning}

The singular values are the square roots of the eigenvalues of
$A^\top A$, and their product is $|\det A|$:

\begin{lstlisting}
np.sqrt(np.linalg.eigvalsh(A.T @ A))   # [1.3028, 2.3028]: ascending
S.prod(), np.linalg.det(A)             # 3.0, 3.0
\end{lstlisting}

For a matrix that is not square, \texttt{full\_matrices=False} returns the
compact form, which keeps only the first $\min(m, n)$ singular vectors on
each side. These are the only ones the product uses:

\begin{lstlisting}
rng = np.random.default_rng(0)
L = rng.standard_normal((100, 3)) @ rng.standard_normal((3, 50))  # rank 3
M = L + 0.1 * rng.standard_normal((100, 50))   # plus small noise

U, S, Vt = np.linalg.svd(M)
U.shape, S.shape, Vt.shape     # (100, 100), (50,), (50, 50)

U, S, Vt = np.linalg.svd(M, full_matrices=False)
U.shape, S.shape, Vt.shape     # (100, 50), (50,), (50, 50)

S[:5]                          # [79.2, 66.3, 63.8, 1.6, 1.6]
np.linalg.matrix_rank(M)       # 50: noise makes every sigma nonzero
\end{lstlisting}

Finally, the truncated SVD. The error matches the Eckart--Young formula,
and with $k = 3$ the approximation is closer to the noise-free matrix
\texttt{L} than \texttt{M} itself is:

\begin{lstlisting}
def truncate(U, S, Vt, k):
    return U[:, :k] * S[:k] @ Vt[:k]       # first k rank-one terms

def rel_err(X, Y):
    return np.linalg.norm(X - Y) / np.linalg.norm(Y)  # Frobenius

for k in [1, 2, 3, 10]:
    Mk = truncate(U, S, Vt, k)
    print(k, rel_err(Mk, M), rel_err(Mk, L))
# 1   0.759  0.758
# 2   0.528  0.526
# 3   0.055  0.018     three terms capture the patterns
# 10  0.045  0.037     extra terms add back noise

M3 = truncate(U, S, Vt, 3)
np.linalg.norm(M - M3), np.sqrt((S[3:]**2).sum())   # equal
rel_err(M, L)                  # 0.058: the noisy matrix is farther
\end{lstlisting}

For large matrices where only the top few singular values are needed,
computing the full SVD is wasteful. \texttt{scipy.sparse.linalg.svds} and
\texttt{sklearn.utils.extmath.randomized\_svd} compute just the leading
$k$ terms, much faster.

\section{Principal Component Analysis}
\label{sec:ch04-pca}

A dataset with $d$ features lives in $\R^d$, but its points rarely fill
that space evenly. The pixels of a handwritten digit are strongly
correlated with their neighbors, and the features of a house move
together: bigger houses tend to have more bedrooms. Most of the variation
in such data happens along a few directions. \emph{Principal component
analysis} (PCA) finds those directions, so that we can describe each data
point with a few numbers instead of $d$, plot high-dimensional data in two
dimensions, or strip away the directions that carry mostly noise. It is
the eigendecomposition of \Cref{sec:ch04-eigendecomposition} and the
truncated SVD of \Cref{sec:ch04-low-rank}, applied to data.

\subsection{Variance Along a Direction}
\label{sec:ch04-pca-variance}

Let the data be the $N \times d$ matrix $X$ whose rows
$\vect{x}_1^\top, \dots, \vect{x}_N^\top$ are the data points. PCA cares
about how the points spread around their center, not where the center
is, so the first step is always to subtract the mean
\[
  \vect{\mu} = \frac{1}{N} \sum_{i=1}^N \vect{x}_i
\]
from every row. Write $\tilde{\vect{x}}_i = \vect{x}_i - \vect{\mu}$ for the
centered points and $\tilde{X}$ for the centered data matrix, whose
columns now each have mean zero.

Pick a unit vector $\vect{w}$. The \emph{score} of point $i$ along
$\vect{w}$ is $z_i = \vect{w}^\top \tilde{\vect{x}}_i$, the signed length
of its projection onto the line through $\vect{w}$
(\Cref{sec:ch03-geometry}). By the row view, all $N$ scores together are
the vector $\vect{z} = \tilde{X}\vect{w}$. They have mean zero, and their
variance, the average squared score, measures how spread out the data is
along $\vect{w}$:
\begin{equation}
  \operatorname{Var}(\vect{z})
  = \frac{1}{N - 1} \sum_{i=1}^N z_i^2
  = \frac{1}{N - 1} \|\tilde{X}\vect{w}\|^2
  = \vect{w}^\top \Bigl( \frac{\tilde{X}^\top \tilde{X}}{N - 1} \Bigr)
    \vect{w} .
  \label{eq:ch04-var-direction}
\end{equation}
Dividing by $N - 1$ rather than $N$ is the usual convention for the
variance of a sample, and the reason for it is a topic of
\Cref{chap:ch06}. It changes every variance by the same factor, so it has
no effect on which directions PCA finds.

The matrix in the middle of \cref{eq:ch04-var-direction} is the
\emph{covariance matrix} of the data,
\[
  C = \frac{1}{N - 1}\, \tilde{X}^\top \tilde{X} .
\]
It is $d \times d$ and symmetric. By the entry-by-entry view of the
product, its diagonal entry $c_{jj}$ is the variance of feature $j$, and
the off-diagonal entry $c_{jk}$ is the \emph{covariance} of features $j$
and $k$: positive when the two tend to be above their means together,
negative when one tends to be high as the other is low.
Equation~\eqref{eq:ch04-var-direction} says that the covariance matrix
knows the spread of the data in every direction, not just along the
coordinate axes: the variance along $\vect{w}$ is $\vect{w}^\top C
\vect{w}$.

\begin{example}
  \label{ex:ch04-pca-data}
  Take the five points $(7, 5)$, $(3, 1)$, $(6, 2)$, $(4, 4)$, and
  $(5, 3)$. Their mean is $\vect{\mu} = (5, 3)$, and the centered points are
  $(2, 2)$, $(-2, -2)$, $(1, -1)$, $(-1, 1)$, and $(0, 0)$. Then
  \[
    \tilde{X}^\top \tilde{X}
    = \begin{bmatrix} 10 & 6 \\ 6 & 10 \end{bmatrix},
    \qquad
    C = \frac14 \begin{bmatrix} 10 & 6 \\ 6 & 10 \end{bmatrix}
      = \begin{bmatrix} 2.5 & 1.5 \\ 1.5 & 2.5 \end{bmatrix}.
  \]
  Both features have variance $2.5$, and their covariance $1.5$ is
  positive: the points lean along the diagonal. Along
  $\vect{e}_1$ the variance is $\vect{e}_1^\top C \vect{e}_1 = 2.5$; along
  the diagonal direction $\vect{w} = \frac{1}{\sqrt 2}(1, 1)$ it is
  $\vect{w}^\top C \vect{w} = \frac12 (2.5 + 1.5 + 1.5 + 2.5) = 4$.
\end{example}

\subsection{Definition}
\label{sec:ch04-pca-def}

The \emph{first principal component} is the direction along which the data
varies the most: the unit vector $\vect{w}$ that maximizes $\vect{w}^\top C
\vect{w}$. Because $C$ is symmetric, the spectral theorem answers this
question directly. Let $\lambda_1 \ge \dots \ge \lambda_d$ be the
eigenvalues of $C$, with orthonormal eigenvectors $\vect{q}_1, \dots,
\vect{q}_d$. Any unit vector can be written as $\vect{w} = c_1 \vect{q}_1 +
\dots + c_d \vect{q}_d$ with $c_1^2 + \dots + c_d^2 = 1$, and then
\[
  \vect{w}^\top C \vect{w}
  = \lambda_1 c_1^2 + \dots + \lambda_d c_d^2
  \le \lambda_1 (c_1^2 + \dots + c_d^2) = \lambda_1 .
\]
The bound is reached at $\vect{w} = \vect{q}_1$. So the direction of
greatest variance is the top eigenvector of the covariance matrix, and the
variance along it is the top eigenvalue. Every eigenvalue of $C$ is a
variance, so none is negative.

Repeating the argument among the directions perpendicular to $\vect{q}_1$
gives $\vect{q}_2$, and so on.

\begin{definition}[Principal components]
  \label{def:pca}
  The \emph{principal components} of a dataset are the unit eigenvectors
  $\vect{q}_1, \dots, \vect{q}_d$ of its covariance matrix $C$, ordered so
  that their eigenvalues satisfy $\lambda_1 \ge \dots \ge \lambda_d$. The
  eigenvalue $\lambda_j$ is the variance of the data along $\vect{q}_j$,
  and the fraction
  \[
    \frac{\lambda_j}{\lambda_1 + \dots + \lambda_d}
  \]
  is the \emph{proportion of variance explained} by component $j$.
\end{definition}

The denominator is the total variance of the data, the sum of the
variances of all the features: $\lambda_1 + \dots + \lambda_d =
\operatorname{tr} C$, by the trace rule of
\Cref{sec:ch04-eigen-compute}. The principal components share that total
out among $d$ perpendicular directions, as unevenly as possible.

To \emph{reduce} the data to $k$ dimensions, put the first $k$ principal
components into the columns of a $d \times k$ matrix $W_k$ and compute the
scores
\[
  \underbrace{Z}_{N \times k}
  = \underbrace{\tilde{X}}_{N \times d} \; \underbrace{W_k}_{d \times k} .
\]
Row $i$ of $Z$ gives the $k$ coordinates of point $i$ along the principal
components. To go back to the original features, map each point to its
projection, $\hat{\vect{x}}_i = \vect{\mu} + W_k \vect{z}_i$, or for all
points at once $\hat{X} = Z W_k^\top + \mathbf{1}\vect{\mu}^\top$.

\begin{example}
  \label{ex:ch04-pca}
  The covariance matrix $C = \bigl[\begin{smallmatrix} 2.5 & 1.5 \\ 1.5 &
  2.5 \end{smallmatrix}\bigr]$ of \Cref{ex:ch04-pca-data} has eigenvalues
  $\lambda_1 = 2.5 + 1.5 = 4$ and
  $\lambda_2 = 2.5 - 1.5 = 1$ along $\vect{q}_1 = \frac{1}{\sqrt2}(1, 1)$
  and $\vect{q}_2 = \frac{1}{\sqrt2}(-1, 1)$, the same directions as the
  matrix of \Cref{sec:ch04-eigen}. The first component explains
  $4 / 5 = 80\%$ of the variance. The scores along $\vect{q}_1$ are
  \[
    \vect{z} = \tilde{X}\vect{q}_1
    = (2\sqrt2,\ -2\sqrt2,\ 0,\ 0,\ 0)
    \approx (2.83,\ -2.83,\ 0,\ 0,\ 0),
  \]
  whose variance is $(8 + 8)/4 = 4 = \lambda_1$, as promised. Keeping
  only this component replaces each point by its projection
  $\vect{\mu} + z_i \vect{q}_1$: the first two points stay where they are,
  and the other three all move to the mean $(5, 3)$. The squared distance
  lost, $2 + 2 + 0 = 4$, is $(N - 1)\lambda_2$, the variance of the
  discarded component.
\end{example}

\paragraph{PCA through the SVD.}
The covariance matrix is $\tilde{X}^\top \tilde{X}$ up to a constant, and
\Cref{sec:ch04-svd-def} showed that the eigenvectors of $\tilde{X}^\top
\tilde{X}$ are the right singular vectors of $\tilde{X}$. So if $\tilde{X} =
U \Sigma V^\top$,
\begin{itemize}
  \item the principal components are the right singular vectors
    $\vect{v}_j$ of the centered data;
  \item the variances are $\lambda_j = \sigma_j^2 / (N - 1)$;
  \item the scores are $Z = \tilde{X} V_k = U_k \Sigma_k$, the first $k$
    columns of $U$ scaled by the singular values.
\end{itemize}
This is how PCA is computed in practice. It never forms $C$, which
squares the singular values and can lose accuracy, and it works with a
truncated SVD when only a few components are needed. It also brings in
the Eckart--Young theorem: the reconstruction $\hat{X}$ from $k$ components
is the best rank-$k$ approximation of the centered data, so no other
choice of $k$ directions loses less.

\subsection{Geometric Interpretation}
\label{sec:ch04-pca-geometry}

\Cref{fig:ch04-pca} shows a centered cloud of $60$ points. Its principal
components are the axes of the elliptical cloud: $\vect{q}_1$ along the
long axis, $\vect{q}_2$ along the short one. The arrows have length
$2\sqrt{\lambda_j}$, twice the standard deviation in each direction, and
the dashed ellipse with these half-lengths outlines where most points
fall. This is the ellipse of \Cref{fig:ch04-eigen} again. The covariance
matrix, like any symmetric matrix, describes an ellipse with its
eigenvectors as axes, and here that ellipse is the shape of the data.

\begin{figure}[ht]
  \centering
  \begin{tikzpicture}[scale=0.85, >=stealth]
    % (a) cloud with principal axes
    \begin{scope}
      \draw[gray!60] (-3.6,0) -- (3.6,0);
      \draw[gray!60] (0,-3.2) -- (0,2.6);
      \draw[dashed, gray!60!black, rotate=32.33] (0,0)
        ellipse [x radius=2.812, y radius=0.817];
      \foreach \p in {
          (0.31,0.69), (0.82,-0.24), (1.17,0.98), (-0.86,-0.12), (0.46,0.48), (-0.07,0.32),
          (-0.95,-0.60), (-0.79,-0.07), (0.16,-0.04), (-0.99,-0.68), (0.11,-0.06), (1.57,1.53),
          (-3.25,-2.93), (-0.10,-0.26), (0.27,0.32), (3.24,1.27), (-1.00,0.64), (0.76,0.86),
          (-0.27,-1.07), (0.24,0.24), (-1.50,-1.22), (0.17,-0.41), (-0.13,0.02), (0.21,-0.13),
          (0.63,0.92), (0.68,-0.04), (1.17,0.43), (1.52,0.30), (1.30,0.78), (-1.62,-1.08),
          (0.04,0.22), (-1.05,-1.21), (0.43,0.02), (0.17,0.58), (-2.32,-1.15), (1.80,0.91),
          (-1.28,-0.27), (0.16,0.65), (-0.08,-0.86), (-0.01,-0.22), (1.06,0.76), (-1.93,-1.76),
          (1.09,1.06), (-0.85,-0.45), (0.53,0.62), (1.18,0.87), (-0.03,-0.13), (2.06,-0.07),
          (-0.17,-0.04), (-2.03,-0.94), (-1.09,-0.09), (-0.31,0.25), (1.63,1.20), (-0.80,-1.30),
          (2.49,1.41), (-0.96,-0.43), (-0.45,0.27), (0.08,0.09), (-1.08,-0.31), (-1.57,-0.48)} {
        \fill[blue!55!black, opacity=0.75] \p circle (1.6pt);
      }
      \draw[->, very thick, orange!85!black] (0,0) -- (2.376,1.504)
        node[above right, font=\small] {$2\sqrt{\lambda_1}\, \vect{q}_1$};
      \draw[->, very thick, green!45!black] (0,0) -- (-0.437,0.691)
        node[above left, font=\small] {$2\sqrt{\lambda_2}\, \vect{q}_2$};
      \node[font=\small] at (0,-3.7) {(a) principal axes};
    \end{scope}
    % (b) projection onto the first component
    \begin{scope}[xshift=8.2cm]
      \draw[gray!60] (-3.6,0) -- (3.6,0);
      \draw[gray!60] (0,-3.2) -- (0,2.6);
      \draw[thick, orange!85!black] (-3.38,-2.14) -- (3.38,2.14)
        node[above, font=\small] {$\vect{q}_1$};
      \foreach \x/\y/\a/\b in {
          0.31/0.69/0.53/0.34, 0.82/-0.24/0.47/0.30, 1.17/0.98/1.28/0.81, -0.86/-0.12/-0.67/-0.42,
          0.46/0.48/0.55/0.35, -0.07/0.32/0.10/0.06, -0.95/-0.60/-0.95/-0.60, -0.79/-0.07/-0.59/-0.37,
          0.16/-0.04/0.10/0.06, -0.99/-0.68/-1.01/-0.64, 0.11/-0.06/0.06/0.03, 1.57/1.53/1.81/1.15,
          -3.25/-2.93/-3.65/-2.31, -0.10/-0.26/-0.19/-0.12, 0.27/0.32/0.34/0.22, 3.24/1.27/2.89/1.83,
          -1.00/0.64/-0.43/-0.27, 0.76/0.86/0.93/0.59, -0.27/-1.07/-0.67/-0.43, 0.24/0.24/0.28/0.18,
          -1.50/-1.22/-1.62/-1.03, 0.17/-0.41/-0.06/-0.04, -0.13/0.02/-0.08/-0.05, 0.21/-0.13/0.09/0.06,
          0.63/0.92/0.87/0.55, 0.68/-0.04/0.47/0.30, 1.17/0.43/1.03/0.65, 1.52/0.30/1.22/0.77,
          1.30/0.78/1.28/0.81, -1.62/-1.08/-1.64/-1.04, 0.04/0.22/0.13/0.08, -1.05/-1.21/-1.30/-0.82,
          0.43/0.02/0.31/0.20, 0.17/0.58/0.38/0.24, -2.32/-1.15/-2.17/-1.38, 1.80/0.91/1.70/1.07,
          -1.28/-0.27/-1.03/-0.65, 0.16/0.65/0.41/0.26, -0.08/-0.86/-0.44/-0.28, -0.01/-0.22/-0.11/-0.07,
          1.06/0.76/1.10/0.70, -1.93/-1.76/-2.17/-1.38, 1.09/1.06/1.25/0.79, -0.85/-0.45/-0.82/-0.52,
          0.53/0.62/0.66/0.42, 1.18/0.87/1.24/0.78, -0.03/-0.13/-0.08/-0.05, 2.06/-0.07/1.44/0.91,
          -0.17/-0.04/-0.14/-0.09, -2.03/-0.94/-1.87/-1.18, -1.09/-0.09/-0.82/-0.52, -0.31/0.25/-0.11/-0.07,
          1.63/1.20/1.70/1.08, -0.80/-1.30/-1.16/-0.73, 2.49/1.41/2.42/1.53, -0.96/-0.43/-0.88/-0.56,
          -0.45/0.27/-0.20/-0.12, 0.08/0.09/0.10/0.06, -1.08/-0.31/-0.91/-0.57, -1.57/-0.48/-1.34/-0.85} {
        \draw[gray!65] (\x,\y) -- (\a,\b);
        \fill[blue!55!black, opacity=0.75] (\x,\y) circle (1.6pt);
        \fill[orange!85!black] (\a,\b) circle (1.3pt);
      }
      \node[font=\small] at (0,-3.7) {(b) projection onto $\vect{q}_1$};
    \end{scope}
  \end{tikzpicture}
  \caption{PCA of $60$ centered points. (a) The principal components are
    the axes of the cloud; each arrow has length twice the standard
    deviation along it, and the dashed ellipse has the same half-lengths.
    (b) Reducing to one dimension replaces each point (blue) by its
    projection onto the first component (orange). PCA chooses the line
    that makes the projections as spread out as possible, which is the
    same as making the gray segments as short as possible.}
  \label{fig:ch04-pca}
\end{figure}

\paragraph{Most variance, least error.}
Panel (b) shows two ways to describe the best line. For each centered
point, the projection onto $\vect{w}$ and the leftover perpendicular piece
form a right triangle, so by Pythagoras
\[
  \|\tilde{\vect{x}}_i\|^2
  = \underbrace{(\vect{w}^\top \tilde{\vect{x}}_i)^2}_{\text{projection}}
  + \underbrace{\|\tilde{\vect{x}}_i - (\vect{w}^\top \tilde{\vect{x}}_i)
    \vect{w}\|^2}_{\text{squared distance to the line}} .
\]
The left side does not depend on $\vect{w}$. Summed over all points, a line
that makes the projections large must make the distances small. Maximizing
the variance of the projections and minimizing the total squared distance
from the points to the line are the same problem, and PCA solves both.
The second view is often the more useful one: PCA finds the
$k$-dimensional flat surface through the mean that fits the data best.

Notice that the distances are measured perpendicular to the line. This
is not the line that linear regression fits, which minimizes vertical
distances because it predicts one variable from the others. PCA treats all
features alike and predicts none.

\paragraph{A rotation of the coordinates.}
Keeping all $d$ components, $Z = \tilde{X}Q$ is just the data written in
new coordinates: since $Q$ is orthogonal, this rotates (or reflects) the
cloud so that its principal axes line up with the coordinate axes. In
the new coordinates the features are \emph{uncorrelated}. The covariance
matrix of the scores is
\[
  \frac{1}{N - 1} Z^\top Z = Q^\top C\, Q = \Lambda ,
\]
which is diagonal: the variances $\lambda_j$ on the diagonal and zero
covariance between any two components. Dimensionality reduction then
simply drops the coordinates with the smallest variance.

\paragraph{What PCA cannot see.}
The geometric picture also shows PCA's limits.
\begin{itemize}
  \item \emph{It depends on units.} Variance is measured in the squared
    units of each feature. In the house data of \Cref{sec:ch03-matrices},
    areas in square meters vary far more than bedroom counts, so the
    first component is almost exactly the area axis and explains $98.5\%$
    of the variance. When features are measured in different units, divide
    each by its standard deviation first, so that every feature has
    variance $1$. Then the first component of the house data weighs all
    three features about equally and explains $92.5\%$.
  \item \emph{It finds flat structure only.} Points spread around a
    circle or along a curved surface have no single best line or plane,
    and PCA will report several components of similar size.
  \item \emph{Large variance is not the same as importance.} PCA keeps the
    directions in which the data varies most, which need not be the ones
    that matter for a given prediction task.
\end{itemize}

\begin{aibox}
  PCA is the first thing to try when looking at high-dimensional data. To
  see how a model organizes its embeddings, project them onto their first
  two principal components and plot them; clusters and directions such as
  ``singular to plural'' can show up immediately. Researchers studying
  the internals of neural networks run PCA on hidden activations for the
  same reason, and nonlinear tools such as t-SNE and UMAP are usually run
  on the top $50$ or so principal components rather than on the raw
  data. Before deep learning, PCA was also a standard preprocessing step,
  used to shrink inputs and remove noise before training a model.
\end{aibox}

\subsection{In Code: PCA from Scratch and with scikit-learn}
\label{sec:ch04-pca-code}

PCA takes a few lines on top of the SVD. The function returns the mean,
the first $k$ principal components as the columns of \texttt{W}, and the
variance along every component:

\begin{lstlisting}
def pca(X, k):
    mu = X.mean(axis=0)
    Xc = X - mu                                   # center the data
    U, S, Vt = np.linalg.svd(Xc, full_matrices=False)
    W = Vt[:k].T                                  # components as columns
    var = S**2 / (len(X) - 1)                     # variance per component
    return mu, W, var
\end{lstlisting}

On the five points of \Cref{ex:ch04-pca-data}:

\begin{lstlisting}
X = np.array([[7.0, 5.0], [3.0, 1.0], [6.0, 2.0],
              [4.0, 4.0], [5.0, 3.0]])

mu, W, var = pca(X, k=1)
mu                         # [5., 3.]
W                          # [[-0.7071], [-0.7071]]: q1, up to sign
var                        # [4., 1.]
var / var.sum()            # [0.8, 0.2]: proportion of variance

Z = (X - mu) @ W           # scores: [[-2.83], [2.83], [0], [0], [0]]
X_hat = Z @ W.T + mu       # [[7, 5], [3, 1], [5, 3], [5, 3], [5, 3]]
((X - X_hat)**2).sum()     # 4.0: squared distance lost
\end{lstlisting}

The SVD has returned $-\vect{q}_1$, so the scores have the opposite sign
from \Cref{ex:ch04-pca}; the projections \texttt{X\_hat} are unchanged.
Going through the covariance matrix gives the same variances, in
ascending order:

\begin{lstlisting}
C = np.cov(X, rowvar=False)        # rows are observations
C                                  # [[2.5, 1.5], [1.5, 2.5]]
np.linalg.eigvalsh(C)              # [1., 4.]
\end{lstlisting}

\begin{note}
  \texttt{np.cov} assumes by default that each \emph{row} is a variable
  and each column an observation, the opposite of the data-matrix
  convention used in this book and in scikit-learn. Pass
  \texttt{rowvar=False}, or you get an $N \times N$ matrix instead of a
  $d \times d$ one.
\end{note}

\paragraph{A real dataset.}
The digits dataset that ships with scikit-learn has $1797$ images of
handwritten digits, each $8 \times 8$ pixels, flattened into vectors in
$\R^{64}$ as in \Cref{sec:ch03-examples}. How many directions does the
data really use?

\begin{lstlisting}
from sklearn.datasets import load_digits

X = load_digits().data         # shape (1797, 64), pixel values 0 to 16
mu, W, var = pca(X, k=2)

ratio = var / var.sum()
ratio[:5]                      # [0.149, 0.136, 0.118, 0.084, 0.058]
np.cumsum(ratio)[[1, 9, 19, 29]]   # [0.285, 0.738, 0.894, 0.959]
np.searchsorted(np.cumsum(ratio), 0.9) + 1    # 21 components for 90%

Z = (X - mu) @ W               # shape (1797, 2): each image as 2 numbers
\end{lstlisting}

Twenty-one of the $64$ directions hold $90\%$ of the variance. Two hold
only $28.5\%$, but even they separate the digits surprisingly well.
\Cref{fig:ch04-digits} plots $300$ of the images by their two scores.
Zeros, fours, and sixes form their own regions, and the overlapping groups
in the middle are digits that look alike in a blurry $8 \times 8$
picture.

\begin{figure}[ht]
  \centering
  \colorlet{digit0}{blue!70!black}
  \colorlet{digit1}{orange!85!black}
  \colorlet{digit2}{green!50!black}
  \colorlet{digit3}{red!75!black}
  \colorlet{digit4}{violet!80!black}
  \colorlet{digit5}{brown!80!black}
  \colorlet{digit6}{magenta!75!black}
  \colorlet{digit7}{gray!60!black}
  \colorlet{digit8}{olive!85!black}
  \colorlet{digit9}{cyan!60!black}
  \begin{tikzpicture}[
    digitlabel/.style={font=\scriptsize\bfseries\sffamily, inner sep=0pt}]
    \begin{axis}[
      width=0.72\textwidth, height=0.6\textwidth,
      xmin=-34, xmax=34, ymin=-32, ymax=30,
      xlabel={first principal component}, ylabel={second principal component},
      tick label style={font=\footnotesize},
      label style={font=\small},
      axis lines=box,
    ]
      \foreach \x/\y/\d in {
      -27.4/-3.8/4, 0.3/10.6/1, 17.2/-4.2/3, -10.3/23.6/1, 3.0/14.8/7, 13.1/-6.9/9,
      -1.3/14.2/7, 8.7/-4.6/5, -4.2/-5.1/8, -8.3/5.3/1, -20.8/7.1/4, -17.3/-13.8/6,
      -9.9/23.7/1, -1.7/12.0/5, 1.7/-22.7/0, 9.4/6.7/2, 7.2/-24.0/0, -17.5/-12.6/6,
      -30.0/-8.0/4, 2.6/17.7/7, -18.4/15.0/4, 1.3/-1.0/5, 0.2/17.5/7, -6.3/11.3/7,
      8.9/11.8/3, 12.9/11.1/2, 14.3/5.1/2, -3.8/-0.8/5, -28.2/-7.8/4, 6.8/3.5/8,
      7.4/20.6/7, 2.6/-7.2/1, -20.8/-7.7/6, -8.5/1.4/5, -17.2/-9.6/6, 6.4/3.8/8,
      11.5/4.3/2, 4.9/1.4/1, 1.9/7.3/5, 6.5/-12.6/9, -0.7/15.0/7, -6.3/7.3/2,
      -25.2/-11.2/4, 15.4/12.4/2, 17.7/-5.2/3, -7.9/5.8/5, 15.9/1.4/3, -0.6/9.4/8,
      -7.2/9.9/1, 8.7/-14.1/9, -5.7/10.4/7, 2.3/3.4/1, 0.3/-19.4/0, 2.2/13.3/7,
      0.9/-22.7/0, 1.4/19.3/7, 23.2/6.5/3, -7.5/18.9/1, 1.3/14.5/7, -28.9/0.7/4,
      -3.8/2.9/8, -3.4/12.0/8, -9.6/-10.6/0, 16.7/-14.0/5, 18.9/-12.7/9, 10.6/4.3/2,
      -14.5/-13.9/6, 15.9/-15.0/5, -2.3/-18.2/0, -22.3/1.3/4, -8.7/13.3/7, -2.8/7.5/2,
      -1.3/-3.2/5, -22.4/6.7/4, -20.5/-4.8/4, -10.9/9.2/1, -6.8/-20.9/0, -14.3/22.4/1,
      9.5/8.6/2, -0.7/26.3/1, 21.9/-2.3/3, -21.3/-7.0/6, -19.0/-18.4/6, -10.2/18.5/1,
      -5.0/-20.8/0, -14.6/-11.4/6, -0.1/18.6/8, -7.7/8.7/5, 16.8/9.0/3, -4.8/1.6/5,
      13.8/6.6/3, -24.1/0.5/4, 5.4/-2.0/3, 7.1/-0.2/8, 12.4/8.6/2, -19.9/-8.2/6,
      -7.6/10.9/1, 5.2/-2.5/1, -18.4/-13.5/6, 5.7/22.4/7, -14.4/-3.2/4, -1.7/8.3/7,
      2.0/10.5/2, 19.5/-3.3/3, 12.8/-10.3/3, 8.6/-18.1/0, -8.1/-8.7/5, -14.1/-8.1/6,
      -16.0/9.9/4, -2.5/10.5/1, 2.0/-6.3/5, 7.9/13.3/9, 11.7/-6.7/3, -7.9/8.4/5,
      2.2/18.7/7, 8.1/21.4/7, 2.8/17.5/7, -14.8/17.5/1, -29.5/-6.8/4, -1.7/12.9/5,
      -9.9/6.1/5, 1.3/-23.8/0, 10.1/5.6/2, 21.3/-1.8/3, 6.5/4.7/2, -4.9/18.1/8,
      -1.7/13.3/8, -16.2/-10.2/6, -2.5/5.4/8, -20.2/6.1/4, -10.1/9.5/8, -17.3/1.5/4,
      6.7/18.2/7, 4.3/2.7/2, 8.1/-4.4/8, 3.9/-2.9/2, -4.1/-18.3/0, 3.8/8.2/8,
      5.7/-8.7/5, 7.7/19.2/9, 2.4/4.9/2, 3.9/10.7/5, -6.6/13.9/2, 10.2/-14.7/9,
      -18.9/-6.0/6, -4.8/11.7/7, 1.2/21.3/7, 7.3/-6.2/1, -3.4/-11.4/5, -15.1/4.3/4,
      13.9/-1.6/3, 20.1/-6.6/3, 8.3/-0.7/8, -14.4/-12.4/6, 14.6/2.6/2, -1.5/14.1/7,
      12.3/-5.3/5, -6.5/15.9/1, 7.1/-2.6/8, 16.7/11.0/2, 18.0/-13.7/9, -13.5/-12.8/6,
      1.6/0.2/1, -5.0/4.7/1, 7.6/-10.5/1, 18.1/2.9/3, 6.1/5.6/2, -14.6/7.9/1,
      -10.4/6.8/8, 19.4/-11.2/9, 10.2/2.4/8, -29.5/-9.4/4, -0.8/17.4/7, -14.1/-16.9/6,
      -14.7/-18.7/6, 12.6/11.9/2, 2.7/-22.6/0, 5.9/-1.2/8, 10.8/6.6/2, 19.1/-10.2/9,
      13.9/-11.9/5, -2.6/-19.2/0, 3.5/11.7/8, 0.8/9.8/8, 1.9/7.0/1, -6.6/-23.5/0,
      -15.6/-11.5/6, -2.6/-20.9/0, -3.9/-23.2/0, -6.7/-23.7/0, 3.1/13.9/7, -15.6/-6.4/6,
      5.8/-0.2/5, 4.1/-8.4/9, -19.8/-11.1/6, -1.7/-18.9/0, 20.9/-3.1/3, -2.6/8.6/7,
      8.2/-4.5/3, 9.0/-12.7/9, -17.4/-15.2/0, -19.4/-8.0/6, -23.3/10.4/4, -14.6/-9.4/0,
      -4.8/5.2/7, -15.8/-17.2/6, -2.4/12.7/7, 3.5/-0.5/5, 23.3/12.9/3, -17.5/-4.6/4,
      4.1/-0.8/9, -1.5/-1.4/2, -6.7/-22.4/0, -1.1/23.6/7, -21.0/-13.3/4, -4.9/-19.3/0,
      -2.5/12.9/7, 6.0/2.9/2, 5.3/-1.2/8, -4.5/11.0/5, 11.4/-5.2/8, 2.7/6.9/2,
      -4.9/20.4/7, -16.0/-15.2/6, -17.0/-17.2/6, 17.7/-14.6/9, -24.1/-6.1/6, 19.7/0.1/3,
      -1.0/-2.6/9, 14.0/-5.2/9, 19.8/7.9/3, -2.6/25.3/1, -13.7/-6.7/4, 14.2/-5.8/3,
      -12.0/-9.0/6, 1.8/-21.6/0, 27.9/-5.6/3, -5.5/9.5/7, 1.4/-22.2/0, -21.7/-7.2/6,
      1.2/17.0/4, -0.5/-3.3/9, -25.1/4.1/4, 8.7/7.7/2, 0.9/-23.8/0, -1.3/5.8/1,
      -4.1/16.2/7, -16.8/-8.0/6, 8.4/-6.2/9, -1.4/2.3/5, -5.5/-13.6/6, 14.1/-11.2/9,
      4.5/21.5/5, 4.3/17.6/5, -3.4/4.8/5, 9.8/-1.4/2, -1.9/18.4/1, -14.7/-15.0/6,
      9.2/-11.3/9, -18.8/-2.4/6, 0.9/-11.5/0, -8.9/21.1/7, -27.3/0.6/4, -4.3/15.6/8,
      -3.9/19.6/7, 3.4/8.8/8, 1.3/-3.8/8, 21.0/-10.0/9, 2.5/4.0/8, 0.2/3.7/8,
      -17.9/-15.6/6, -1.2/5.3/8, 0.8/9.2/1, 6.8/-10.7/1, -5.5/26.8/1, 9.9/4.9/2,
      2.6/2.6/5, 3.7/2.3/2, -1.8/10.8/8, 3.2/11.3/9, 4.6/-23.0/0, 5.2/17.6/8,
      -6.7/-12.1/5, 12.1/10.5/2, 3.7/15.0/8, -2.7/9.7/8, 22.2/-1.7/3, -13.6/-12.1/6,
      19.4/-9.7/9, 24.8/-4.8/3, 2.5/-24.7/0, -24.4/2.9/4, 18.1/-12.1/9, -3.6/11.3/8,
      -23.3/-4.3/4, -14.6/-17.1/6, -0.3/-20.0/0, -9.4/-0.9/6, -3.8/0.2/5, 6.0/11.2/5
      } {
        \edef\temp{\noexpand\node[digitlabel, text=digit\d]
          at (axis cs:\x,\y) {\d};}
        \temp
      }
    \end{axis}
  \end{tikzpicture}
  \caption{$300$ handwritten digits from scikit-learn's digits dataset,
    each $64$-dimensional image placed at its scores on the first two
    principal components and drawn as its label. The two directions
    explain only $28.5\%$ of the variance, yet several digits already
    occupy their own regions.}
  \label{fig:ch04-digits}
\end{figure}

Reconstructing from $k$ components and measuring the error relative to
the spread of the data shows the trade-off. By Eckart--Young, the error is
exactly $\sqrt{1 - \text{(proportion of variance kept)}}$:

\begin{lstlisting}
for k in [2, 10, 21, 30]:
    mu, W, var = pca(X, k)
    X_hat = (X - mu) @ W @ W.T + mu
    err = np.linalg.norm(X - X_hat) / np.linalg.norm(X - mu)
    print(k, err)
# 2   0.846
# 10  0.512
# 21  0.311
# 30  0.202
\end{lstlisting}

In practice, use scikit-learn's \texttt{PCA}, which wraps the same SVD,
picks consistent signs, and can choose $k$ for you from a target
proportion of variance:

\begin{lstlisting}
from sklearn.decomposition import PCA

pca2 = PCA(n_components=2).fit(X)
pca2.explained_variance_           # [179.0, 163.7]: our var[:2]
pca2.explained_variance_ratio_     # [0.149, 0.136]
pca2.components_.shape             # (2, 64): components as ROWS
Z = pca2.transform(X)              # centers, then projects

PCA(n_components=0.9).fit(X).n_components_   # 21
\end{lstlisting}

Note that \texttt{components\_} stores the principal components as rows,
like \texttt{Vt}, so it is the transpose of our \texttt{W}. When features
are measured in different units, standardize them first, for example with
\texttt{StandardScaler} from \texttt{sklearn.preprocessing}, as discussed in
\Cref{sec:ch04-pca-geometry}.

\section{Summary}
\label{sec:ch04-summary}

\begin{itemize}
  \item An \emph{eigenvector} of a square matrix is a nonzero vector whose
    direction the matrix does not change: $A\vect{v} = \lambda\vect{v}$.
    The \emph{eigenvalue} $\lambda$ says whether vectors along it are
    stretched, shrunk, flipped, or crushed to zero.
  \item Repeated multiplication acts on eigenvectors as powers,
    $A^k\vect{v} = \lambda^k\vect{v}$, so the eigenvalue of largest size
    dominates over many steps. This drives power iteration, the exploding
    and vanishing gradients of recurrent networks, and the step-size limit
    of gradient descent, which is set by the largest eigenvalue of the
    matrix of second derivatives.
  \item Eigenvalues solve $\det(A - \lambda I) = 0$. They add up to the
    trace and multiply to the determinant. In practice, compute them with
    \texttt{np.linalg.eigh} for symmetric matrices, which returns them in
    increasing order, and \texttt{np.linalg.eig} otherwise.
  \item A general matrix may have too few eigenvectors, or none, but a
    \emph{symmetric} matrix has $n$ orthogonal ones (the \emph{spectral
    theorem}). It factors as $A = Q\Lambda Q^\top$: rotate, stretch along
    the axes, rotate back.
  \item Every matrix, of any shape, has a \emph{singular value
    decomposition} $A = U\Sigma V^\top$, with orthogonal $U$ and $V$ and
    singular values $\sigma_1 \ge \sigma_2 \ge \dots \ge 0$. It pairs input
    directions with output directions, $A\vect{v}_i = \sigma_i\vect{u}_i$.
    The $\vect{v}_i$ are eigenvectors of $A^\top A$, with eigenvalues
    $\sigma_i^2$.
  \item Geometrically, every matrix is a rotation, a stretch along the
    axes, and another rotation. It maps the unit sphere to an ellipsoid
    whose half-axes are $\sigma_i \vect{u}_i$; $\sigma_1$ is the largest
    stretch, and the number of nonzero singular values is the
    \emph{rank}.
  \item The SVD writes a matrix as a sum of rank-one pieces
    $\sigma_i \vect{u}_i \vect{v}_i^\top$. Keeping the $k$ largest gives the
    best rank-$k$ approximation (Eckart--Young), with an error made of the
    discarded singular values. This underlies compression, denoising, and
    LoRA.
  \item \emph{Principal component analysis} centers the data and finds the
    directions of greatest variance: the eigenvectors of the covariance
    matrix $C = \tilde{X}^\top\tilde{X}/(N - 1)$, or equivalently the right
    singular vectors of the centered data. Eigenvalue $\lambda_j$ is the
    variance along component $j$.
  \item Projecting onto the first $k$ components keeps as much variance as
    possible and, equivalently, minimizes the squared perpendicular
    distance from the points to a $k$-dimensional flat surface. The new
    coordinates are uncorrelated.
  \item PCA depends on the units of the features, so standardize features
    measured in different units. It finds only flat structure, and large
    variance need not mean importance. In code, mind that
    \texttt{np.linalg.svd} returns $V^\top$, and that both it and
    scikit-learn's \texttt{PCA} store components as rows.
\end{itemize}
